% Chapter 20: The Sato-Tate Conjecture and the Langlands Program
\let\LocalMacro\relax

\chapter{Sato-Tate Conjecture and the Langlands Program}

\section{The Problem}

\subsection{Informal}
Imagine a doughnut-shaped curve defined by a simple equation. For each prime number, you can count how many points sit on that curve when you do arithmetic with that prime---think of it like a clock with a prime number of hours. The count is always close to a predictable value, but it wobbles up and down in a seemingly random way. The question is: \emph{as you range over larger and larger primes, do those wobbles follow a predictable pattern?}

The Sato--Tate conjecture, proposed in 1960, says yes---for a typical curve, the wobbles follow a specific bell-shaped pattern, much like human heights cluster around an average. This simple-sounding claim took over fifty years to prove and became a landmark test for the Langlands program, one of the deepest unifying visions in mathematics.

\subsection{Formal}
Let $E/\mathbb{Q}$ be an elliptic curve without complex multiplication. For each prime $p$ of good reduction, let $a_p(E) = p + 1 - \#E(\mathbb{F}_p)$. The \emph{Sato-Tate conjecture} (1960) predicts the distribution of the normalized Frobenius traces:

\begin{equation}
\frac{a_p(E)}{2\sqrt{p}} \in [-1, 1]
\end{equation}

\begin{conjecture}[Sato-Tate]
For a non-CM elliptic curve $E/\mathbb{Q}$, the sequence $\left\{ \frac{a_p(E)}{2\sqrt{p}} \right\}_{p \text{ good}}$ is equidistributed in $[-1, 1]$ with respect to the \emph{Sato-Tate measure}
\[
\mu_{\text{ST}} = \frac{2}{\pi} \sqrt{1 - t^2} \, dt.
\]
\end{conjecture}

\subsection{Historical Context}
Sato and Tate independently observed this distribution numerically for $y^2 = x^3 - x$ (which has CM) and for non-CM curves like $y^2 = x^3 - 4x^2 + 2x$. The conjecture remained open for 50 years, becoming a cornerstone test for the Langlands program.

\subsection{What Was Known}
The conjecture was proposed by Sato and Tate independently in 1960. By the 2000s, it had been proved for many specific curves using modular form techniques. But a general proof---one that works for all typical curves---required the full machinery of the Langlands program, including modularity theorems and potential modularity arguments that were still being developed.

\subsection{Why It Matters}
The Sato--Tate conjecture is a test case for the Langlands program, one of the deepest unifying visions in mathematics. Proving it required developing new techniques that have since been applied to many other problems in number theory. The result also has practical implications: understanding the distribution of points on curves is essential for cryptography and coding theory.

\subsection{Simplified Version}
For elliptic curves with complex multiplication (CM curves), like $y^2 = x^3 - x$, the distribution was proved by Serre in 1968---the Frobenius traces follow a specific arcsine law. The open case was \emph{non-CM} curves, where the distribution follows the Sato--Tate density $\frac{2}{\pi}\sin^2\theta$. Taylor, Harris, Shepherd-Barron, and Clozel proved this in 2008--2011 using automorphic forms and potential automorphy.

\section{Mathematical Theory}


\subsection{Prerequisites}
This chapter assumes familiarity with:
\begin{itemize}
\item Algebraic number theory (Galois representations, L-functions)
\item Elliptic curves and modular forms
\item Basic representation theory of compact groups
\item The classical Langlands program (reciprocity laws)
\end{itemize}

\subsection{Galois Representations and Automorphic Forms}
For an elliptic curve $E/\mathbb{Q}$, the $p$-adic Tate module $T_p(E)$ gives a Galois representation
\[
\rho_{E,p}: \text{Gal}(\overline{\mathbb{Q}}/\mathbb{Q}) \to \text{GL}_2(\mathbb{Z}_p).
\]
The Sato-Tate conjecture is equivalent to the statement that the restriction of $\rho_{E,p}$ to the decomposition group at $p$ matches a modular form of weight 2.

\subsection{The Role of the Langlands Program}
The Langlands program predicts a correspondence:
\begin{itemize}
\item \textbf{Galois side}: $n$-dimensional Galois representations $\rho: \text{Gal}(\overline{\mathbb{Q}}/\mathbb{Q}) \to \text{GL}_n(\mathbb{C})$.
\item \textbf{Automorphic side}: Automorphic representations $\pi$ of $\text{GL}_n(\mathbb{A}_{\mathbb{Q}})$.
\end{itemize}
For $n=2$, the Sato-Tate conjecture is a special case of the \emph{modularity theorem} (formerly Taniyama-Shimura-Weil) plus the \emph{generalized Sato-Tate conjecture} for higher-genus curves.

\subsection{Equidistribution and Moments}
The Sato-Tate measure $\mu_{\text{ST}}$ has moments:
\[
\int_{-1}^1 t^k \, d\mu_{\text{ST}}(t) = \begin{cases}
0 & k \text{ odd} \\
C_{k/2} & k \text{ even}
\end{cases}
\]
where $C_n = \frac{1}{n+1}\binom{2n}{n}$ are Catalan numbers. Verifying these moments for $a_p(E)$ is equivalent to the conjecture.

\section{The Solution}

\subsection{What Was Missing Before}

The Sato-Tate conjecture asked whether the normalized Frobenius traces $a_p(E)/(2\sqrt{p})$ of a non-CM elliptic curve $E/\mathbb{Q}$ are equidistributed in $[-1,1]$ with respect to the measure $\frac{2}{\pi}\sqrt{1-t^2}\,dt$. Before the proof, the situation was fragmented. Serre had proved the conjecture for CM curves in 1968 \cite{serre1968}: for curves with complex multiplication (like $y^2 = x^3 - x$), the Frobenius traces follow an arcsine law, not the Sato-Tate distribution. For individual non-CM curves, modular form techniques could verify equidistribution computationally and sometimes theoretically. But nobody could prove the conjecture for \emph{all} non-CM curves simultaneously. The gap was: the equidistribution statement was known in special cases but completely out of reach in general.

\subsection{Why Existing Methods Failed}

The core obstacle was that standard modular form techniques work curve by curve. For a specific elliptic curve $E$, the modularity theorem (Wiles 1995, BCDT 2001) tells you that $E$ corresponds to a weight-2 modular form $f$, and the equidistribution of $a_p(E)$ follows from the equidistribution of Hecke eigenvalues of $f$---a consequence of the spectral theorem for Hecke operators on $L^2$ of modular forms. But this argument requires you to know that $E$ is modular \emph{and} that the relevant automorphic representation has the right properties. For a single curve, Wiles's modularity theorem suffices. But to prove Sato-Tate for \emph{every} non-CM curve, you need a uniform argument that does not depend on knowing the specific modular form associated to $E$. The problem was that the modularity of $E/\mathbb{Q}$ does not automatically give you the equidistribution result you need: you need the Galois representation $\rho_{E,\lambda}$ to be automorphic in a way that the trace formula can handle, and for general (non-semistable) curves over general number fields, this was unknown.

\subsection{What Was Novel}

The proof by Clozel, Harris, Shepherd-Barron, and Taylor (2008--2011) introduced a powerful new strategy: \textbf{potential automorphy}. The idea is deceptively simple. Instead of proving that $\rho_{E,\lambda}$ is automorphic over $\mathbb{Q}$ directly, you show that it \emph{becomes} automorphic after a solvable base change to a larger number field $F$. More precisely:

\begin{enumerate}
\item \textbf{Potential automorphy} \cite{taylor2002}: Taylor proved in 2002 that for any elliptic curve $E/\mathbb{Q}$ and any prime $\lambda$, the Galois representation $\rho_{E,\lambda}$ restricted to $\text{Gal}(\overline{\mathbb{Q}}/F)$ is potentially automorphic for some totally real field $F$. This means there exists a solvable extension $F/\mathbb{Q}$ such that $\rho_{E,\lambda}|_F$ is automorphic.
\item \textbf{Base change}: Automorphy is preserved under solvable base change (via Langlands's base change for $\text{GL}_2$). So you can transfer the automorphic representation from $F$ back down to $\mathbb{Q}$.
\item \textbf{Stable trace formula}: Arthur's stable trace formula for $\text{GL}_2$ over CM fields lets you compare the automorphic representation on $F$ with one coming from a CM field, where Sato-Tate is already known.
\item \textbf{Equidistribution transfer}: On the automorphic side, equidistribution of Hecke eigenvalues follows from the spectral theorem. The trace formula lets you transfer this equidistribution back to the Galois side, proving Sato-Tate for $E$.
\end{enumerate}

The key insight is that you do not need to prove automorphy directly---you only need it \emph{potentially}, i.e., after a base change. This sidesteps the hardest part of the Langlands correspondence and lets you use the trace formula as a bridge between what you know (CM cases) and what you want (all non-CM cases).

\subsection{Student-Friendly Explanation}

Think of it this way. You want to prove that the ``wobble'' in the number of points on an elliptic curve, as you vary the prime, follows a specific bell-shaped pattern. The natural approach is: use the modularity theorem to match each elliptic curve to a modular form, then use harmonic analysis on modular forms to prove the equidistribution. This works for individual curves, but to prove it for \emph{all} curves at once, you need a trick.

The trick is \emph{potential automorphy}. Imagine you have a puzzle that is hard to solve in its original setting, but if you move it to a bigger room (a larger number field), it becomes easy. Potential automorphy says: after moving to a suitable extension field, the Galois representation ``looks like'' an automorphic form, and on the automorphic side, the equidistribution follows from basic harmonic analysis (the eigenvalues of Hecke operators are equidistributed by the spectral theorem). The stable trace formula is the elevator that carries you back from the bigger room to the original one.

So the logical chain is: (1) the Galois representation becomes automorphic after a base change; (2) on the automorphic side, equidistribution is automatic; (3) the trace formula transfers this equidistribution back to the Galois side. The CM case (proved by Serre in 1968) provides the starting point: you know equidistribution holds for CM curves, and the trace formula lets you bootstrap from there to all non-CM curves.

\subsection{Modularity and Potential Modularity}
\begin{itemize}
\item \textbf{Wiles \cite{wiles1995}}: Modularity for semistable elliptic curves (proved FLT).
\item \textbf{Breuil-Conrad-Diamond-Taylor \cite{bcdt2001}}: Full modularity for all elliptic curves over $\mathbb{Q}$.
\item \textbf{Taylor (2002)}: Potential modularity for elliptic curves over totally real fields.
\end{itemize}

\subsection{The Proof (2008-2011)}
The Sato-Tate conjecture was proved in a series of papers by \textbf{Clozel, Harris, Shepherd-Barron, and Taylor} (2008, 2011).

\begin{theorem}[Sato-Tate \cite{barnetlamb2011}]
For any non-CM elliptic curve $E/\mathbb{Q}$, the Sato-Tate conjecture holds.
\end{theorem}

The proof strategy:
\begin{enumerate}
\item \textbf{Potential automorphy}: Show that the Galois representation $\rho_{E,\lambda}$ is potentially automorphic (becomes automorphic after a solvable base change).
\item \textbf{Stable trace formula}: Use Arthur's stable trace formula for $\text{GL}_2$ to relate Galois representations to automorphic forms on $\text{GL}_2$ over CM fields.
\item \textbf{Equidistribution}: Transfer the equidistribution of Frobenius traces from the automorphic side (where it follows from the spectral theorem) to the Galois side.
\item \textbf{CM case}: The CM case was already known by Serre (1968) and uses Hecke characters.
\end{enumerate}

\begin{highlightbox}
The proof is a landmark of the Langlands program: it uses the full power of the stable trace formula, base change, and potential automorphy to transfer a Galois-theoretic question to an automorphic one where harmonic analysis applies.
\end{highlightbox}

\subsection{Generalized Sato-Tate}
The conjecture generalizes to higher-genus curves and higher-dimensional abelian varieties. The equidistribution measure is determined by the \emph{Sato-Tate group}, a compact subgroup of $\text{USp}(2g)$.

\section{Impact}
\begin{itemize}
\item \textbf{Langlands program}: First major application of the stable trace formula to a classical number-theoretic problem.
\item \textbf{Modularity}: Extended modularity techniques to non-semistable curves and higher-dimensional varieties.
\item \textbf{Computational number theory}: Verified Sato-Tate distributions for millions of curves.
\item \textbf{Arithmetic statistics}: Led to precise conjectures on the distribution of $a_p$ for higher-genus curves (FKRS conjecture).
\end{itemize}

\section{Exercises}

\begin{exercise}[Basic]
For $E: y^2 = x^3 - x$ (CM by $\mathbb{Z}[i]$), compute the first 10 values of $a_p$ and explain why they do \emph{not} follow the Sato-Tate distribution.
\end{exercise}

\begin{exercise}[Intermediate]
Show that the Sato-Tate measure $\frac{2}{\pi}\sqrt{1-t^2} dt$ is the eigenvalue distribution of a random matrix in $\text{SU}(2)$ with Haar measure.
\end{exercise}

\begin{exercise}[Challenge]
Explain why the stable trace formula is needed instead of the ordinary trace formula to prove Sato-Tate.
\end{exercise}

\section{Timeline}

\begin{tabularx}{\linewidth}{lX}
\toprule
\textbf{Year} & \textbf{Event} \\ \midrule
1960 & Sato and Tate independently conjecture the distribution \\ 
1968 & Serre proves Sato-Tate for CM elliptic curves \\ 
1995 & Wiles proves modularity for semistable curves \\ 
2001 & BCDT prove modularity for all elliptic curves over $\mathbb{Q}$ \\ 
2008 & Clozel-Harris-Shepherd-Barron-Taylor announce Sato-Tate proof \\ 
2011 & Full proof published in JAMS \\ 
2010s & Generalized Sato-Tate for higher-genus curves (FKRS) \\ \bottomrule
\end{tabularx}

\section{Citations}

\begin{enumerate}
\item Clozel, L., Harris, M., \& Taylor, R. (2008). ``Automorphy for some $\ell$-adic lifts of automorphic mod $\ell$ Galois representations.'' Publications Mathématiques de l'IHÉS, 108(1), 1--181.
\item Barnet-Lamb, T., Geraghty, D., Harris, M., \& Taylor, R. (2011). ``A family of Calabi-Yau varieties and potential automorphy II.'' Publications Mathématiques de l'IHÉS, 116(1), 119--180.
\item Serre, J.-P. (1968). ``Abelian $\ell$-adic representations and elliptic curves.'' W. A. Benjamin.
\item Harris, M. (2011). ``The Sato-Tate conjecture.'' Notices of the AMS, 58(8), 1046--1056.
\end{enumerate}
